\chapter{TypeScript}

JavaScript ha un sistema di tipi peculiare, che abbiamo imparato ad
\sout{amare} accettare.

È \term{tipizzato dinamicamente}:

\begin{affianco}[0.5][c]
\begin{lstlisting}[style=js]
let x;
console.log(typeof(x));   // undefined
x = 1.25e7;
console.log(typeof(x));   // number
x = "Hello Web Technologies!";
console.log(typeof(x));   // string
\end{lstlisting}
\risultato
\begin{immagine}[La stessa variabile cambia tipo nel tempo.]
\begin{tikzpicture}[
    t/.style={rounded corners=2pt, draw=primary!60!black, fill=primary!12!white,
      line width=0.7pt, font=\FiraCodeNL\scriptsize, minimum width=1.7cm, minimum height=0.5cm},
    a/.style={-{Stealth[length=4pt,width=3.5pt]}, line width=0.7pt, draw=neutral!45!black}]
  \node[t] (a) at (0,0) {undefined};
  \node[t] (b) at (2.3,0) {number};
  \node[t] (c) at (4.6,0) {string};
  \draw[a] (a) -- (b); \draw[a] (b) -- (c);
  \node[mnota, anchor=north] at (a.south) {\code{let x}};
  \node[mnota, anchor=north] at (b.south) {\code{x = 1.25e7}};
  \node[mnota, anchor=north] at (c.south) {\code{x = "Hello\ldots"}};
\end{tikzpicture}
\end{immagine}
\end{affianco}

Ed è anche \term{tipizzato debolmente}, con una spiccata predilezione per le
conversioni implicite:

\begin{affianco}[0.6][c]
\begin{lstlisting}[style=js]
let y = "100" - 1 + "42";
// string - number -> number, poi number + string -> string
console.log(`${y}: ${typeof(y)}`);   // 9942: string
\end{lstlisting}
\risultato
\begin{immagine}[Due conversioni implicite in una riga.]
\begin{tikzpicture}[
    v/.style={mcod},
    a/.style={-{Stealth[length=4pt,width=3.5pt]}, line width=0.7pt, draw=neutral!45!black},
    l/.style={mnota, above, text=accent!70!black}]
  \node[v] (a) at (0,0) {"100" - 1};
  \node[v] (b) at (2.2,0) {99 + "42"};
  \node[v] (c) at (4.3,0) {"9942"};
  \draw[a] (a) -- node[l]{\(-\) numerico} (b);
  \draw[a] (b) -- node[l]{\(+\) concatena} (c);
\end{tikzpicture}
\end{immagine}
\end{affianco}

\section{Il problema}

Queste caratteristiche rendono il linguaggio pratico per piccole manipolazioni
di pagine web: funziona e basta, senza bisogno di essere prolissi dichiarando
tipi o conversioni esplicite.

Man mano che la complessità dei programmi cresce, però:

\begin{itemize}
  \item è più facile introdurre bug insidiosi;
  \item l'esperienza di sviluppo peggiora: poco supporto dall'IDE, pochissimi
        errori individuati prima dell'esecuzione;
  \item il codice diventa meno manutenibile e più difficile da capire.
\end{itemize}

\begin{esempio}{Bug che si scoprono solo a runtime}
\begin{affianco}[0.6][c]
\codicepiccolo
\begin{lstlisting}[style=js, name=bugs.js]
let message = "Hello Web Technologies!";
console.log(message.toLowerCase());   // hello web technologies!
message();       // TypeError: message is not a function (a runtime)

let employee = { name: "Jordan Belfort", role: "Stockbroker" };
employee.nome = "The Wolf of Wall Street";   // nessun errore!
console.log(employee.name);                  // Jordan Belfort

let team = ["Jordan Belfort", "Donnie Azoff", "Chester Ming"];
team.add("Nicky Koskoff");   // TypeError: team.add is not a function
\end{lstlisting}
\risultato
\begin{console}
hello web technologies!
Uncaught TypeError: message is
  not a function
\end{console}
\wtnota{L'esecuzione si ferma al primo errore;\\
        il refuso \code{nome} non ne produce affatto.}
\end{affianco}

Il secondo caso è il più insidioso: un refuso in un nome di proprietà
(\code{nome} invece di \code{name}) non produce \textbf{alcun} errore. Sarebbe
bello intercettare questi bug \textbf{prima} dell'esecuzione.
\end{esempio}

\section{I transpiler}

JavaScript ha cominciato a essere usato per programmi sempre più complessi, e
gli sviluppatori hanno sentito il bisogno di superare questi limiti. Sostituire
JavaScript con un linguaggio nuovo non era praticabile: era già supportato
ovunque.

Sono così nati molti linguaggi che \textbf{compilano verso JavaScript
standard}: CoffeeScript, TypeScript, Dart, Elm e circa altri 350.

\begin{figure}[H]
  \centering
  \begin{tikzpicture}[font=\Lato\small,
      n/.style={wtnode, minimum width=3.4cm, minimum height=1.3cm},
      a/.style={-{Stealth[length=6pt,width=5pt]}, line width=1pt,
                draw=neutral!45!black}]
    \node[n] (src) at (0,0)
      {codice in un nuovo\\ linguaggio, con\\ funzionalità in più};
    \node[n, draw=orange-gradient-6, fill=orange-gradient-1!40!white]
      (tr) at (5.0,0) {\textbf{Transpiler}};
    \node[n] (js) at (10.0,0)
      {codice JavaScript\\ (gira ovunque giri\\ JavaScript)};
    \draw[a] (src) -- (tr);
    \draw[a] (tr) -- (js);
  \end{tikzpicture}
  \caption{Un transpiler traduce da un linguaggio a un altro, non a codice
           macchina.}
  \label{fig:transpiler}
\end{figure}

\dfn{TypeScript}{
  \term{TypeScript} nasce nel 2012, progettato e sviluppato da Microsoft (lo
  sviluppatore principale è Anders Hejlsberg). È un \textbf{superset} di
  JavaScript con tipizzazione statica, interfacce, enumerazioni e generici, e
  \textbf{compila verso JavaScript}. È di gran lunga la variante di JavaScript
  più diffusa.}

Essendo un superset, \textbf{qualunque codice JavaScript valido è anche codice
TypeScript valido}: si può adottare gradualmente, un file alla volta.

\section{Primi passi}

TypeScript si installa come pacchetto npm:

\begin{lstlisting}[style=shell]
npm install -g typescript
\end{lstlisting}

Il comando rende disponibile globalmente il compilatore \code{tsc}. I file
TypeScript hanno tipicamente estensione \code{.ts}.

\begin{esempio}{Il compilatore trova un bug nel JavaScript puro}
\begin{affianco}[0.48]
\begin{lstlisting}[style=ts, name=hello.ts]
function greet(entity, date) {
  console.log(`Hello ${entity}, today is ${date}!`);
}

greet("Web Technologies");
\end{lstlisting}
\risultato
\begin{console}[Terminale]
$ node hello.ts
Hello Web Technologies, today is undefined!
\end{console}
\begin{console}[Terminale]
$ tsc hello.ts
hello.ts:5:1 - error TS2554:
  Expected 2 arguments, but got 1.

5 greet("Web Technologies");
  ~~~~~~~~~~~~~~~~~~~~~~~~~
  hello.ts:1:24
    An argument for 'date' was
    not provided.
\end{console}
\end{affianco}

A destra le due prove: eseguito con Node, il file funziona (è JavaScript
valido) ma stampa un risultato sbagliato; compilato con \code{tsc}, l'errore
viene segnalato prima dell'esecuzione.


Il compilatore ha individuato il bug \textbf{pur avendo scritto soltanto
JavaScript standard}, senza una sola annotazione di tipo.
\end{esempio}

\subsection{Downleveling}

Corretto il bug e compilato il file, \code{tsc} produce un \code{hello.js}
eseguibile con Node. Il codice generato conserva commenti e indentazione, e
sembra scritto da una persona. Un dettaglio però cambia:

\begin{affianco}[0.42]
\codicepiccolo
\begin{lstlisting}[style=ts, name=hello.ts]
console.log(`Hello ${entity}, today is ${date}!`);
\end{lstlisting}
\risultato
\codicepiccolo
\begin{lstlisting}[style=js, name=hello.js]
console.log("Hello ".concat(entity, ", today is ").concat(date, "!"));
\end{lstlisting}
\end{affianco}

I template literal sono una funzionalità di ES6 e, per impostazione predefinita,
TypeScript ha come bersaglio ES5 (del 2009; nelle versioni precedenti alla 5.0
era addirittura ES3, del 1999). TypeScript sa \textbf{riscrivere}
il codice da versioni più recenti a versioni più vecchie: è il
\term{downleveling}. Possiamo usare le funzionalità moderne e far girare il
risultato anche su browser antichi.

Il bersaglio si può specificare:

\begin{lstlisting}[style=shell]
tsc --target es6 hello.ts
\end{lstlisting}

\section{Dichiarare i tipi}

I tipi primitivi più comuni sono \code{string}, \code{number} e
\code{boolean}.

\begin{lstlisting}[style=ts, name=primitive_types.ts]
let courseName: string = "Web Technologies";
let credits: number = 6;
let isGreat: boolean = true;
\end{lstlisting}

A questo punto il file \textbf{non è più JavaScript valido}: eseguirlo
direttamente con Node produce un \code{SyntaxError}. Va compilato prima.

\begin{affianco}[0.47][c]
\begin{lstlisting}[style=ts]
credits = "9 CFU";   // in JavaScript sarebbe lecito
\end{lstlisting}
\risultato
\begin{console}[Terminale]
$ tsc primitive_types.ts
primitive_types.ts:5:1 - error TS2322:
  Type 'string' is not assignable to
  type 'number'.
\end{console}
\end{affianco}

\subsection{Array}

Il tipo di un array si indica con la sintassi \code{tipo[]}.

\begin{lstlisting}[style=ts, name=array.ts]
let langs: string[] = ["TypeScript", "JavaScript", "PHP", "Java"];
let primes: number[] = [2, 5, 7, 97, 829, 839];

langs.push({ name: "C#" });
// error TS2345: Argument of type '{ name: string; }' is not assignable
// to parameter of type 'string'.
\end{lstlisting}

\subsection{Il tipo any}

TypeScript include un tipo speciale: \code{any}. Quando un valore è di tipo
\code{any} possiamo farci qualunque cosa e il compilatore non protesterà.

\begin{lstlisting}[style=ts]
let x: any;
x = { name: "Angular", lang: "TypeScript" };
x.foo();              // compila, ma solleva TypeError a runtime

x = "Hello!";
let n: number = x;    // compila: stiamo mettendo una stringa in un number
\end{lstlisting}

\warning{\code{any} disattiva i controlli}{
  Usare \code{any} equivale a rinunciare a tutto ciò per cui si è adottato
  TypeScript. È talvolta inevitabile (per esempio interfacciandosi con librerie
  non tipizzate), ma va considerato un'eccezione consapevole, non una scorciatoia
  di comodo.}

\subsection{Funzioni}

Si possono specificare i tipi sia degli ingressi sia del valore restituito.

\begin{lstlisting}[style=ts]
function greet(name: string): string {
  return `Hello ${name.toUpperCase()}`;
}

console.log(greet(42));
// error TS2345: Argument of type 'number' is not assignable to
// parameter of type 'string'
\end{lstlisting}

\subsection{Inferenza automatica}

Quando non si fornisce un'annotazione esplicita, TypeScript prova a
\textbf{dedurre} il tipo dal contesto.

\begin{affianco}[0.62][c]
\codicepiccolo
\begin{lstlisting}[style=ts]
function greet(name) {         // function greet(name: any): string
  return `Hello ${name}`;
}

let arr = [1, 2, 3, 4, 5];     // let arr: number[]
let course = "Web Technologies";   // let course: string

course = { name: "Web Technologies", credits: 6 };
// TS2322: Type '{name: string; credits: number;}' is not assignable
// to type 'string'
\end{lstlisting}
\risultato
\begin{immagine}[Ciò che mostra l'editor passando sopra\\ ai nomi: i tipi dedotti.]
\begin{tikzpicture}[
    tip/.style={fill=consolebg, text=consolefg, rounded corners=2pt,
      font=\FiraCodeNL\scriptsize, inner sep=3pt, anchor=west}]
  \node[tip] at (0,0)    {function greet(name: any): string};
  \node[tip] at (0,-0.6) {let arr: number[]};
  \node[tip] at (0,-1.2) {let course: string};
\end{tikzpicture}
\end{immagine}
\end{affianco}

Quando non si specifica un tipo e TypeScript non riesce a dedurlo, il
compilatore ripiega di norma su \code{any}. Il flag \code{noImplicitAny}
permette di segnalare come errore ogni \code{any} implicito.

\begin{console}[Terminale]
$ tsc --noImplicitAny implicit_any.ts
implicit_any.ts:3:16 - error TS7006: Parameter 'name' implicitly has an
'any' type.
\end{console}

\section{Tipi oggetto}

Un tipo oggetto si definisce elencando le sue proprietà e i loro tipi.

\begin{affianco}[0.6][c]
\codicepiccolo
\begin{lstlisting}[style=ts, name=object_types.ts]
function printCourse(course: { name: string, credits: number }) {
  console.log(`Course name: ${course.name}`);
  console.log(`Course credits: ${course.credits}`);
}

printCourse({ name: "Web Technologies", credits: 6 });   // ok

printCourse({ name: "Software Engineering" });
// TS2345: Property 'credits' is missing

printCourse({ name: "Algebra", credits: 6, year: "First" });
// TS2353: Object literal may only specify known properties
\end{lstlisting}
\risultato
\begin{immagine}[L'oggetto deve avere esattamente la forma attesa.]
\begin{tikzpicture}[
    f/.style={rounded corners=2pt, draw=neutral!55!black, fill=white, line width=0.6pt,
      font=\FiraCodeNL\scriptsize, inner sep=3pt, anchor=west}]
  \node[f, draw=primary!60!black, fill=primary!10!white] (t) at (0,0) {\{ name: string, credits: number \}};
  \node[mnota, anchor=south west] at (t.north west) {atteso};
  \node[f] (a) at (0,-0.75) {\{ name, credits \}};   \node[anchor=west] at ([xshift=3pt]a.east) {\mok};
  \node[f] (b) at (0,-1.35) {\{ name \}};            \node[anchor=west] (bx) at ([xshift=3pt]b.east) {\merr};
  \node[mnota, anchor=west] at (bx.east) {manca \code{credits}};
  \node[f] (c) at (0,-1.95) {\{ name, credits, year \}}; \node[anchor=west] (cx) at ([xshift=3pt]c.east) {\merr};
  \node[mnota, anchor=west] at (cx.east) {\code{year} in più};
\end{tikzpicture}
\end{immagine}
\end{affianco}

Una proprietà si può dichiarare \textbf{opzionale} aggiungendo un \code{?} dopo
il nome.

\begin{lstlisting}[style=ts]
function printCourse(course: { name: string, credits?: number }) {
  console.log(`Course name: ${course.name}`);
  console.log(`Course credits: ${course.credits}`);
  // andrebbe verificato se course.credits e' undefined!
}

printCourse({ name: "Software Engineering" });   // ok: stampa "undefined"
\end{lstlisting}

\section{Tipi unione}

I \term{tipi unione} si ottengono combinando due o più tipi, elencati e separati
da \code{|}. Rappresentano valori che possono appartenere a uno qualunque dei
tipi combinati.

\begin{lstlisting}[style=ts, name=union.ts]
function printError(code: number | string) {
  console.log(`Error code: ${code}`);
}

printError(404);                 // ok
printError("401 Unauthorized");  // ok
printError({ code: "403" });     // errore di compilazione
\end{lstlisting}

TypeScript garantisce che le operazioni su un tipo unione siano valide per
\textbf{ciascuno} dei tipi possibili:

\begin{lstlisting}[style=ts]
function printError(code: number | string) {
  console.log(`Error code: ${code.toUpperCase()}`);
  // TS2339: Property 'toUpperCase' does not exist on type 'string | number'
}
\end{lstlisting}

In questi casi bisogna \term{restringere} il tipo con controlli espliciti (per
esempio \code{typeof} o \code{Array.isArray()}); TypeScript sa dedurre tipi
specifici nei singoli rami.

\begin{affianco}[0.6][c]
\codicepiccolo
\begin{lstlisting}[style=ts]
function printError(code: number | string) {
  if (typeof code === "string") {
    console.log(`Error: ${code.toUpperCase()}`);   // qui code: string
  } else {
    console.log(`Error: ${code}`);                 // qui code: number
  }
}
\end{lstlisting}
\risultato
\begin{immagine}[Dentro ogni ramo il tipo è più preciso.]
\begin{tikzpicture}[
    t/.style={rounded corners=2pt, draw=primary!60!black, fill=primary!10!white,
      line width=0.7pt, font=\FiraCodeNL\scriptsize, inner sep=3pt},
    a/.style={-{Stealth[length=4pt,width=3.5pt]}, line width=0.7pt, draw=neutral!45!black}]
  \node[t] (u) at (0,0) {code: number | string};
  \node[t, fill=green-gradient-1!45!white, draw=green-gradient-6] (s) at (-1.3,-1.4) {code: string};
  \node[t, fill=accent!12!white, draw=accent!55!black] (n) at (1.4,-1.4) {code: number};
  \draw[a] (u) -- node[mnota, left, align=right]{\code{typeof}\\ \code{=== "string"}} (s);
  \draw[a] (u) -- node[mnota, right]{altrimenti} (n);
  \node[mnota, anchor=north] at (s.south) {\code{toUpperCase()} \mok};
\end{tikzpicture}
\end{immagine}
\end{affianco}

\section{Alias di tipo e interfacce}

Volendo riusare lo stesso tipo più volte, si può assegnargli un nome con la
parola chiave \code{type}.

\begin{affianco}[0.6][c]
\codicepiccolo
\begin{lstlisting}[style=ts, name=aliases.ts]
type Credits = number | string;

type Course = {
  name: string,
  credits: Credits
}

function printCourse(course: Course) {
  console.log(`Name: ${course.name}, Credits: ${course.credits}`);
}

let webtech: Course = { name: "Web Technologies", credits: 6 };
let softeng: Course = { name: "Software Engineering", credits: "10 CFU" };
\end{lstlisting}
\risultato
\begin{immagine}[Un alias può comporre altri alias.]
\begin{tikzpicture}[    f/.style={rounded corners=2pt, draw=primary!60!black, fill=primary!8!white,
      line width=0.7pt, font=\FiraCodeNL\scriptsize, inner sep=3pt, align=left},
    a/.style={-{Stealth[length=4pt,width=3.5pt]}, line width=0.7pt, draw=neutral!45!black}]
  \node[f] (c) at (0,0) {type Course\\ \ \ name: string\\ \ \ credits: \textcolor{accent!70!black}{Credits}};
  \node[f, fill=accent!12!white, draw=accent!55!black] (k) at (3.3,-0.35) {type Credits\\ \ \ number | string};
  \draw[a] ([yshift=-0.35cm]c.east) -- (k.west);
\end{tikzpicture}
\end{immagine}
\end{affianco}

Le \term{interfacce} sono un altro modo di dare un nome ai tipi oggetto.

\begin{lstlisting}[style=ts, name=interfaces.ts]
interface Course {
  name: string;
  credits: number;
}

function printCourse(course: Course) {
  console.log(`Name: ${course.name}, Credits: ${course.credits}`);
}
\end{lstlisting}

\subsection{Tipizzazione strutturale}

\dfn{Tipizzazione strutturale}{
  TypeScript è un linguaggio a \term{tipizzazione strutturale}: per stabilire la
  compatibilità fra tipi guarda soltanto alla \textbf{struttura} e alle
  capacità, non ai \textbf{nomi}. Java, al contrario, ha un sistema di tipi
  \textit{nominativo}.}

\begin{affianco}[0.55][c]
\begin{lstlisting}[style=ts, name=structural.ts]
interface Person { name: string; age: number; }
interface Pet    { name: string; age: number; }

function printPerson(p: Person) {
  console.log(`Name: ${p.name}, age: ${p.age}`);
}

let person: Person = { name: "Janet", age: 20 };
let pet: Pet       = { name: "Chuck", age: 3 };

printPerson(person);   // Name: Janet, age: 20
printPerson(pet);      // Name: Chuck, age: 3   (!)
\end{lstlisting}
\risultato
\begin{immagine}[Stessa forma, quindi stesso tipo:\\ i nomi non contano.]
\begin{tikzpicture}[    f/.style={rounded corners=2pt, draw=primary!60!black, fill=primary!8!white,
      line width=0.7pt, font=\FiraCodeNL\scriptsize, inner sep=3pt, align=left},
    a/.style={-{Stealth[length=4pt,width=3.5pt]}, line width=0.7pt, draw=neutral!45!black}]
  \node[f] (p) at (0,0) {Person\\ \ \ name: string\\ \ \ age: number};
  \node[f, fill=accent!12!white, draw=accent!55!black] (q) at (3.0,0) {Pet\\ \ \ name: string\\ \ \ age: number};
  \node at (1.5,0) {\Large =};
  \node[mnota, anchor=north] at (1.5,-0.8) {TypeScript \mok\quad Java \merr};
\end{tikzpicture}
\end{immagine}
\end{affianco}

Un \code{Pet} viene accettato dove è atteso un \code{Person} perché ne ha la
stessa forma. In Java sarebbe un errore di compilazione.

\subsection{Estendere interfacce e tipi}

\begin{affianco}[0.5][c]
\begin{lstlisting}[style=ts]
interface Pet {
  name: string
}

interface Bird extends Pet {
  flies: boolean
}

let chuck: Bird = { name: "Chuck", flies: true };
let quentin: Bird = { flies: false };
// TS2741: Property 'name' is missing
\end{lstlisting}
\risultato
\begin{immagine}[\code{Bird} ha tutto ciò che ha \code{Pet}, più \code{flies}.]
\begin{tikzpicture}[
    r/.style={draw=primary!60!black, fill=primary!8!white, line width=0.7pt,
      font=\FiraCodeNL\scriptsize, minimum width=2.6cm, minimum height=0.45cm},
    n/.style={r, fill=accent!12!white, draw=accent!55!black}]
  \node[r] (n1) at (0,0) {name: string};
  \node[n, anchor=north] (n2) at ([yshift=0.7pt]n1.south) {flies: boolean};
  \draw[wtbrace] ([xshift=-3pt]n1.north west) -- ([xshift=-3pt]n1.south west)
    node[midway, left=4pt, mnota]{Pet};
  \draw[wtbrace, decoration={mirror=false}] ([xshift=3pt]n1.north east) -- ([xshift=3pt]n2.south east)
    node[midway, right=4pt, mnota]{Bird};
\end{tikzpicture}
\end{immagine}
\end{affianco}

Alias di tipo e interfacce sono molto simili e spesso intercambiabili. Una
distinzione chiave: i \code{type} \textbf{non possono essere riaperti} per
aggiungere proprietà, mentre le interfacce sono sempre estendibili.

\begin{lstlisting}[style=ts]
interface Pet { name: string }
interface Pet { age: number }     // ok: le due dichiarazioni si fondono

type Pet2 = { name: string }
type Pet2 = { age: number }       // TS2300: Duplicate identifier
\end{lstlisting}

Gli alias si estendono con i \term{tipi intersezione}, creati con \code{\&}:

\begin{lstlisting}[style=ts]
type Pet = { name: string }

type Bird = Pet & {
  flies: boolean
}

let chuck: Bird = { name: "Chuck", flies: true };
\end{lstlisting}

\section{Asserzioni di tipo}

Talvolta lo sviluppatore ne sa più del compilatore. Per esempio, usando
\code{document.getElementById()} TypeScript può dedurre soltanto che la chiamata
restituirà un \code{HTMLElement}, mentre noi sappiamo che quell'elemento è di un
tipo più specifico.

\begin{lstlisting}[style=ts]
const nameInput = document.getElementById("name") as HTMLInputElement;
\end{lstlisting}

\warning{Non è un'asserzione nel senso dei test}{
  Il termine è fuorviante: non viene verificato nulla a runtime. È soltanto
  un'indicazione al compilatore, valida a tempo di compilazione. Se ci
  sbagliamo, l'errore si manifesterà comunque durante l'esecuzione.}

\section{Tipi letterali ed enumerazioni}

Oltre ai tipi \code{string} e \code{number}, si può fare riferimento a stringhe
e numeri \textbf{specifici} in posizione di tipo.

\begin{lstlisting}[style=ts]
let alignment: "center";
alignment = "center";
alignment = "left";
// TS2322: Type '"left"' is not assignable to type '"center"'
\end{lstlisting}

Da solo non è molto utile, ma combinato con i tipi unione esprime un concetto
molto potente: tipi che corrispondono a un \textbf{insieme di valori noti}.

\begin{lstlisting}[style=ts]
type Alignment = "center" | "left" | "right";

let alignment: Alignment;
alignment = "left";    // ok
alignment = "right";   // ok
alignment = "centre";
// TS2820: Type '"centre"' is not assignable to type 'Alignment'.
// Did you mean '"center"'?
\end{lstlisting}

Le \term{enumerazioni} permettono di definire un insieme di costanti con un
nome.

\begin{affianco}[0.55][c]
\begin{lstlisting}[style=ts, name=enums.ts]
interface Order { isPremium: boolean }

enum Priority {
  Low,
  Medium,
  High
}

function computePriority(s: Order): Priority {
  if (s.isPremium)
    return Priority.High;
  return Priority.Low;
}

console.log(Priority.Medium);   // 1
\end{lstlisting}
\risultato
\begin{immagine}[Senza valori espliciti, le costanti\\ sono numerate da 0.]
\begin{tikzpicture}
  \mchiave{1.4cm}
  \oggetto{e}{(0,0)}{2.4cm}{enum Priority}{Low/0, Medium/1, High/2}
\end{tikzpicture}
\end{immagine}
\end{affianco}

\section{Tipi funzione}

Come si dichiara che una funzione riceve una \textbf{callback} con certi
parametri e un certo risultato? Con le \term{espressioni di tipo funzione}.

\begin{lstlisting}[style=ts]
type stringDecoratorFunction = (x: string, mode: boolean) => string;

let fn: () => string;                  // funzione senza argomenti, ritorna string
fn = () => { return "Hello Web Tech" };       // ok
fn = (x: string) => { return `Hello ${x}` };  // errore di tipo
\end{lstlisting}

\begin{affianco}[0.62][c]
\codicepiccolo
\begin{lstlisting}[style=ts, name=function_types.ts]
enum Case { Uppercase, Lowercase }

function GREET(name: string, decorator: (x: string, mode: Case) => string) {
  return decorator(name, Case.Uppercase);
}

function stringDecorator(x: string, mode: Case) {
  if (mode === Case.Lowercase)
    return `>>> Hello ${x} <<<`.toLowerCase();
  else
    return `>>> Hello ${x} <<<`.toUpperCase();
}

console.log(GREET("Web Technologies", stringDecorator));
\end{lstlisting}
\risultato
\begin{console}
>>> HELLO WEB TECHNOLOGIES <<<
\end{console}
\wtnota{\code{GREET} chiama la callback\\ con \code{Case.Uppercase}.}
\end{affianco}

\section{Generici}

Buona parte dell'ingegneria del software mira a costruire software
\textbf{riusabile}: sviluppare componenti capaci di operare su dati presenti e
futuri è essenziale per i sistemi di grandi dimensioni. I \term{generici} sono
uno degli strumenti principali per scrivere codice che funzioni con una varietà
di tipi.

\begin{esempio}{La funzione identità}
Supponiamo di dover realizzare una funzione che restituisce ciò che riceve.

\begin{lstlisting}[style=ts]
function identity(x) {        // function identity(x: any): any
  return x;
}

let y: string = "Hello";
let z = identity(y);          // z: any (!)
\end{lstlisting}

Usare \code{any} è abbastanza generico, ma \textbf{perdiamo l'informazione} sul
tipo al ritorno. Dare un tipo specifico risolve il problema ma funzionerebbe
solo con le stringhe.

TypeScript permette di definire \term{variabili di tipo}, dichiarate dopo il
nome della funzione fra \code{<} e \code{>}.

\begin{affianco}[0.55][c]
\codicepiccolo
\begin{lstlisting}[style=ts, name=generics.ts]
function identity<Type>(x: Type): Type {
  return x;
}

let s: string = identity<string>("Hello");   // Type esplicito
let n: number = identity(42);                // dedotto automaticamente
let o: { title: string, artist: string } = identity({
  title: "Sultans of swing",
  artist: "Dire Straits"
});
\end{lstlisting}
\risultato
\begin{immagine}[Il tipo in uscita è quello in ingresso.]
\begin{tikzpicture}[
    t/.style={rounded corners=2pt, draw=primary!60!black, fill=primary!10!white,
      line width=0.7pt, font=\FiraCodeNL\scriptsize, inner sep=3pt},
    fn/.style={mfun},
    a/.style={-{Stealth[length=4pt,width=3.5pt]}, line width=0.7pt, draw=neutral!45!black}]
  \foreach \t/\y in {string/0, number/-0.7, {\{ title, artist \}}/-1.4} {
    \node[t, anchor=east] (i) at (-0.1,\y) {\t};
    \node[t, anchor=west] (o) at (2.5,\y) {\t};
    \draw[a] (i.east) -- (0.35,\y); \draw[a] (2.05,\y) -- (o.west);
  }
  \node[fn, minimum height=2.0cm, minimum width=1.6cm] at (1.2,-0.7) {identity};
\end{tikzpicture}
\end{immagine}
\end{affianco}

Non è la stessa cosa di usare \code{any}: qui l'informazione sul tipo di
ingresso viene \textbf{preservata}.
\end{esempio}

Le variabili di tipo possono essere più di una, e funzionano anche con gli
array:

\begin{lstlisting}[style=ts]
function merge<T1, T2>(x: T1, y: T2): { field: T2, otherField: T1 } {
  return { field: y, otherField: x };
}

let x = merge("Sultans of Swing", 42);
// x ha tipo { field: number, otherField: string }

function getFirst<T>(arr: T[]): T {
  return arr[0];
}

let first: string = getFirst(["hello", "web", "technologies"]);
\end{lstlisting}

\subsection{Classi generiche}

\begin{lstlisting}[style=ts, name=generic_class.ts]
class CustomCollection<T> {
  list: T[] = [];

  addElement(element: T): number {
    return this.list.push(element);
  }

  getRandomElement(): T {
    let selectedIndex: number = Math.floor(Math.random() * this.list.length);
    return this.list[selectedIndex];
  }
}

let c = new CustomCollection<string>();
c.addElement("hello");
c.addElement("web");
console.log(c.getRandomElement());
\end{lstlisting}

\subsection{Vincoli sui generici}

\begin{affianco}[0.55][c]
\begin{lstlisting}[style=ts, name=constraints.ts]
class Animal { species: string; }
class Bird extends Animal  { canFly: boolean; }
class Snake extends Animal { isVenomous: boolean; }

function animalInfo<T extends Animal>(animal: T): void {
  console.log(animal.species);
}

animalInfo(new Bird());              // ok
animalInfo({ species: "Spider" });   // ok anche questo (!)
\end{lstlisting}
\risultato
\begin{immagine}[Basta avere \code{species}: per la tipizzazione\\
                 strutturale anche l'oggetto letterale è un \code{Animal}.]
\begin{tikzpicture}[
    c/.style={rounded corners=2pt, draw=primary!60!black, fill=primary!8!white,
      line width=0.7pt, font=\FiraCodeNL\scriptsize, inner sep=3pt},
    a/.style={-{Stealth[length=5pt,width=4pt,open]}, line width=0.7pt, draw=neutral!45!black}]
  \node[c] (an) at (1.4,0) {Animal\\[-1pt] \scriptsize species};
  \node[c] (b) at (0,-1.4) {Bird};
  \node[c] (s) at (1.4,-1.4) {Snake};
  \node[c, dashed, fill=white] (l) at (3.1,-1.4) {\{ species \}};
  \draw[a] (b) -- (an); \draw[a] (s) -- (an);
  \draw[a, dashed] (l) -- (an);
\end{tikzpicture}
\end{immagine}
\end{affianco}

L'ultima riga funziona per via della tipizzazione strutturale: l'oggetto
letterale ha la stessa forma di \code{Animal}, quindi soddisfa il vincolo.

\section{Altri tipi da conoscere}

\begin{center}
\renewcommand{\arraystretch}{1.45}
\begin{tabular}{@{}l@{\hspace{1.2cm}}>{\raggedright\arraybackslash}p{8.6cm}@{}}
  \textbf{Tipo} & \textbf{Significato} \\
  \code{void}    & tipo di ritorno delle funzioni che non restituiscono un
                   valore \\
  \code{unknown} & rappresenta qualunque valore possibile: simile ad
                   \code{any}, ma \textbf{più sicuro}, perché non consente
                   alcuna operazione senza prima restringere il tipo \\
  \code{never}   & rappresenta valori che non vengono mai osservati: il tipo di
                   una funzione che non ritorna mai, o di un ramo irraggiungibile \\
\end{tabular}
\end{center}

\begin{lstlisting}[style=ts, name=unknown.ts]
function greet(a: any) {
  console.log(a.name);       // compila, ma possibile errore a runtime
}

function saferGreet(a: unknown) {
  console.log(a.name);       // TS18046: 'a' is of type unknown
}

type namedObject = { name: string };

function notSoSafeGreet(a: unknown) {
  console.log((a as namedObject).name);   // asserzione esplicita
}
\end{lstlisting}

\begin{lstlisting}[style=ts, name=never.ts]
function f(x: string | number) {
  if (typeof x === "string") {
    return x.toLowerCase();
  } else if (typeof x === "number") {
    return x.toPrecision(2);
  } else {
    return x;   // qui x ha tipo never: nessun altro caso e' possibile
  }
}

function g(x: any): never {   // g non ritorna mai: solleva sempre un errore
  throw new Error("Oops!");
}
\end{lstlisting}

\section{Livelli di severità}

Utenti diversi si aspettano cose diverse da TypeScript. Per impostazione
predefinita offre un'esperienza \textbf{opzionale}: i tipi non sono obbligatori
e si ricorre ad \code{any} quando non è possibile dedurre un tipo preciso. Il
flag \code{noImplicitAny} rende più severo questo comportamento.

Per non essere troppo invasivo, alcuni controlli sono disattivati di
partenza. Per esempio \code{null} e \code{undefined} possono essere assegnati a
qualunque tipo: dimenticarsi di gestirli esplicitamente è la causa di
innumerevoli bug, tanto che il concetto è stato definito \guillemotleft l'errore da un
miliardo di dollari\guillemotright.

Il flag \code{strictNullChecks} impone che \code{null} e \code{undefined} siano
gestiti in modo esplicito.

\begin{affianco}[0.55][c]
\begin{lstlisting}[style=ts, name=strict.ts]
let character = x.list.find((val: string) => {
  return val === "Amelie"
});

// character potrebbe essere undefined!
console.log(character.toUpperCase());
\end{lstlisting}
\risultato
\begin{console}[Terminale]
$ tsc -target es6 strict.ts
$ tsc -target es6 -strictNullChecks strict.ts
strict.ts:13:13 - error TS18048:
  'character' is possibly 'undefined'.
\end{console}
\end{affianco}

Senza il flag il codice compila e il bug esplode a runtime; con il flag viene
intercettato subito. È il motivo per cui nei progetti nuovi conviene attivare
fin dall'inizio la modalità severa.

\section{Riferimenti}

\begin{itemize}
  \item \textbf{The TypeScript Handbook}. \\
        \urlx{https://www.typescriptlang.org/docs/handbook/intro.html} \\
        Non serve conoscerlo per intero: va usato come riferimento per le parti
        discusse in questo capitolo.
  \item \textbf{Lecture Notes for the Comparative Studies of Programming
        Languages Course}, Paquet e Mokhov (2010). \\
        \urlx{https://arxiv.org/pdf/1007.2123.pdf} \\
        Sezione 3.3 per approfondire i sistemi di tipi (duck typing, sistemi
        strutturali contro nominativi). Non richiesto per il corso.
  \item \textbf{To Type or Not to Type? A Systematic Comparison of the Software
        Quality of JavaScript and TypeScript Applications on GitHub}, Justus
        Bogner e Manuel Merkel (2022). \\
        \urlx{https://dl.acm.org/doi/pdf/10.1145/3524842.3528454}
\end{itemize}
